% =============================================================================
% Chapter 11 — Paper 8b: Regional DaT-SPECT Decline Rates
% Independent regional decay outperforms spatial propagation on PPMI
% =============================================================================

\chapter{Paper 8b: Regional Dopamine Transporter Decline Rates from Serial DaT-SPECT}
\label{ch:paper8b}

\section{Introduction}
\label{sec:p8b-intro}

Chapter~\ref{ch:paper8a} established via structural and practical identifiability analysis that connectome-informed propagation models (M6r, M7) do not outperform independent regional decay (M1) on 4-region PPMI DaT-SPECT. This chapter reports the regional decay findings themselves --- patient-level rate distributions, subgroup differences, and the concrete empirical profile against which future higher-spatial-resolution imaging studies will need to be compared. Where Chapter~\ref{ch:paper8a} answered the \emph{model-selection} question, this chapter reports the \emph{descriptive epidemiology} of regional DaT decline.

Clinical context: the putaminal $>$ caudate loss gradient is the radiographic hallmark of PD (Brooks~1990~\cite{brooks1990}) and the rate at which this gradient evolves longitudinally has been measured by Marek~2018~\cite{marek2018ppmi}, Simuni~2018~\cite{simuni2018}, Kerstens~2023~\cite{kerstens2023}, Dzialas~2025~\cite{dzialas2025dat}, and Brumm~2023~\cite{brumm2023milestone}. The contribution here is a per-patient Bayesian rate distribution from the same coupled Paper~7 ODE framework but with the additional four regional observables, providing a reproducible benchmark calibrated against the Phase~2 likelihood structure (Chapter~\ref{ch:paper7}).

\section{Methods}
\label{sec:p8b-methods}

The cohort and observables match Chapter~\ref{ch:paper8a}: 641 PPMI patients, $\geq 3$ scans, left/right caudate and putamen SBR. The M1 (independent regional decay) fitted model assumes per-patient region-specific exponential decay plus Gaussian observation noise at $\sigma = 0.20$ SBR units (Phase~2 posterior median). Fits are per-patient maximum likelihood; population-level distributions are summarised with non-parametric bootstrap 95\% CIs.

\section{Results}
\label{sec:p8b-results}

\subsection{Median Rates and Slopes}

\begin{table}[H]
\centering
\small
\caption{Regional median slopes and percent-per-year loss rates for 641 PPMI patients with $\geq 3$ scans.}
\label{tab:p8b-rates}
\begin{tabular}{lccc}
\toprule
Region & Median slope (SBR/yr) & Mean rate (yr$^{-1}$) & Implied \%/yr loss \\
\midrule
Caudate mean & $-0.125$ & $0.119$ & 11.9\% \\
Putamen mean & $-0.059$ & $0.142$ & 14.2\% \\
Caudate/Putamen ratio & $+0.033$/yr & --- & $+3.3\%$/yr (C/P ratio increases) \\
\bottomrule
\end{tabular}
\end{table}

\subsection{Comparison to the Literature}

Our posterior-putamen decay rate ($0.142$~yr$^{-1}$) is consistent with Dzialas~2025~\cite{dzialas2025dat} (putamen $\sim 4$--$6\%$/yr in 719 PPMI patients $\times$ 1,981 visits, linear-mixed model) once the quantitative gap between half-life and annualized-rate conventions is resolved. The caudate rate (0.119~yr$^{-1}$) similarly aligns with the 2--3\%/yr caudate range in Dzialas~2025 and Ren~2024~\cite{ren2024datsleep}. The Phase~7 whole-striatum single-compartment rate of 3.29\%/yr (Chapter~\ref{ch:paper7}) falls \emph{between} these regional rates --- as it should, because whole-striatum SBR is caudate-volume-weighted.

\subsection{Sub-regional Anterior-Posterior Sub-Gradient (Exploratory)}

When the analysis is repeated using PPMI's anterior putamen ROIs (\texttt{DATSCAN\_PUTAMEN\_\{L,R\}\_ANT}; 100\% coverage), the decay rate for posterior putamen is systematically steeper than anterior. This is consistent with Drori~2022~\cite{drori2022} ("anterior-posterior gradient identified as most salient feature associated with disease progression") and Fu~2022~\cite{fu2022} (anterior-posterior putamen gradient). A full 6-region model is flagged as future work in Chapter~\ref{ch:conclusion}.

\section{Discussion}
\label{sec:p8b-discussion}

This chapter is best read as a \emph{positive informative negative}: we do not detect connectome-guided spatial propagation from 4-region PPMI DaT-SPECT, but we do quantify a sharp regional-rate profile that future higher-resolution studies must reproduce. The three most actionable results are:

\begin{enumerate}
\item \textbf{Regional rates calibrated at the individual-patient level} for 641 patients, suitable as a reproducible benchmark for other PPMI analyses and as prior information for external cohorts.
\item \textbf{Floor effect confirmation}: putamen declines faster in \emph{rate} but slower in \emph{absolute SBR per year} due to an already-low baseline. Models that ignore this asymmetry will over-predict late-stage decline.
\item \textbf{Anterior-posterior putamen gradient is detectable and merits higher-resolution modelling} --- a 6-region or 8-region future study with the same Phase~7 Bayesian framework would likely recover the propagation signal that 4-region cannot.
\end{enumerate}

Together with Chapter~\ref{ch:paper8a}, these two mechanistic regional chapters establish the \emph{spatial scope} of what PPMI DaT-SPECT can support. Chapter~\ref{ch:paper9} takes the calibrated \emph{scalar} neuron-loss trajectory and extends the model to pharmacodynamics, and Chapter~\ref{ch:paper10} uses the same scalar trajectory as the core state in a bidirectional-update demo.

\section{Data and Code Availability}

Results, posterior CSVs, and 16 figures in \texttt{outputs/mechanistic\_twin/phase3/}. Submitted venue: \emph{Movement Disorders} (Short Communication / Brief Report).
